../cictikz/src/cictikz/data/tex/cictikz_preamble.tex

% The LELO_EX schematic the tutorial builds.
%
% The devices are the JNWATR_NCH_4C5F0 symbol itself, converted from the
% xschem library rather than redrawn, so what a reader places on the
% canvas and what the book prints cannot drift apart. Gate on the left,
% drain and source to the right, bulk between them - tied to the source
% here, which is the short loop beside each device.
%
% The circuit is a current mirror with three instances of that one
% device. The input branch xo1 is a single device and is diode
% connected; the two output instances are two devices each, because the
% bus notation xo0[1:0] creates two. Four output devices to one input
% device is where the 5 uA in and 20 uA out of the port names comes
% from.

\begin{circuitikz}[american, thick, transform shape, circuit ee IEC]

  ../cictikz/src/cictikz/data/tex/cictikz_lib.tex
  % Generated by py/gen_symlib.py from the JNW xschem symbol libraries.
% Do not edit: change the symbol and run `make symlib`.
%
% These are the symbols a reader places in xschem during the tutorial,
% converted rather than redrawn so the book and the tool cannot drift
% apart. Each macro is a path fragment starting at the current point and
% exporting its pins as coordinates: (name_g), (name_d), (name_s),
% (name_b).


% Generated from JNWATR_NCH_4C5F0.sym - edit the symbol, not this file.
\newcommand{\jnwNchFiveF}[1]{
  ++(0,0) coordinate (#1_org)
  (#1_org) ++(0.825,-0.275) -- ++(0.275,-0.1375)
  (#1_org) ++(1.1,-0.4125) -- ++(-0.275,-0.1375)
  (#1_org) ++(0.55,0.4125) -- ++(0.55,0)
  (#1_org) ++(1.1,0.4125) -- ++(0,0.4125)
  (#1_org) ++(0.55,-0.4125) -- ++(0.55,0)
  (#1_org) ++(1.1,-0.4125) -- ++(0,-0.4125)
  (#1_org) ++(0.55,-0) -- ++(0.55,0)
  (#1_org) ++(0.55,-0.4125) -- ++(0,0.825)
  (#1_org) ++(0.4125,0.4125) -- ++(0,-0.825)
  (#1_org) ++(0,-0) -- ++(0.4125,0)
  (#1_org) ++(1.1,0.825) coordinate (#1_d)
  (#1_org) ++(0,0) coordinate (#1_g)
  (#1_org) ++(1.1,-0.825) coordinate (#1_s)
  (#1_org) ++(1.1,0) coordinate (#1_b)
  (#1_org)
}

% Generated from JNWATR_PCH_4C5F0.sym - edit the symbol, not this file.
\newcommand{\jnwPchFiveF}[1]{
  ++(0,0) coordinate (#1_org)
  (#1_org) ++(0.55,-0.4125) -- ++(0.55,0)
  (#1_org) ++(1.1,-0.4125) -- ++(0,-0.4125)
  (#1_org) ++(0.55,0.4125) -- ++(0.55,0)
  (#1_org) ++(1.1,0.4125) -- ++(0,0.4125)
  (#1_org) ++(0.55,-0) -- ++(0.55,0)
  (#1_org) ++(1.0175,0.55) -- ++(-0.275,-0.1375)
  (#1_org) ++(0.7425,0.4125) -- ++(0.275,-0.1375)
  (#1_org) ++(0.55,-0.4125) -- ++(0,0.825)
  (#1_org) ++(0.4125,0.4125) -- ++(0,-0.825)
  (#1_org) ++(0,-0) -- ++(0.1925,0)
  (#1_org) ++(0.3025,-0) circle (0.11)
  (#1_org) ++(1.1,-0.825) coordinate (#1_d)
  (#1_org) ++(0,0) coordinate (#1_g)
  (#1_org) ++(1.1,0.825) coordinate (#1_s)
  (#1_org) ++(1.1,0) coordinate (#1_b)
  (#1_org)
}


  \def\xin{0}       % xo1, the diode connected input
  \def\xoa{3.2}     % xo0[1:0]
  \def\xob{6.4}     % xo2[1:0]
  \def\ygate{-1.9}  % the shared gate rail
  \def\yvss{-2.9}   % VSS
  \def\yout{2.4}    % the output rail

  % one instance: #1 x, #2 node name, #3 instance name, #4 device count
  \newcommand{\instance}[4]{
    \draw (#1,0) \jnwNchFiveF{#2};
    %- gate out to the left, then down to the rail the three share
    \draw (#2_g) -- ({#1-0.5},0) -- ({#1-0.5},\ygate);
    %- source down to VSS
    \draw (#2_s) -- ({#1+1.1},\yvss);
    %- bulk tied to the source, beside the device
    \draw (#2_b) -- ({#1+1.55},0) -- ({#1+1.55},-0.825) -- (#2_s);
    \node[anchor=south east, scale=0.85] at ({#1+1.0},1.0) {#3};
    \node[anchor=north east, scale=0.8, gray!50!black] at ({#1+1.0},-1.05) {#4};
  }

  \instance{\xin}{XO1}{xo1}{1 device}
  \instance{\xoa}{XO0}{xo0[1:0]}{2 devices}
  \instance{\xob}{XO2}{xo2[1:0]}{2 devices}

  % ---- the shared gate rail ----
  \draw ({\xin-0.5},\ygate) -- ({\xob-0.5},\ygate);
  \fill ({\xoa-0.5},\ygate) circle (0.075);

  % ---- VSS ----
  \draw ({\xin-1.3},\yvss) -- ({\xob+1.55},\yvss);
  \fill ({\xin+1.1},\yvss) circle (0.075);
  \fill ({\xoa+1.1},\yvss) circle (0.075);
  \fill ({\xob+1.1},\yvss) circle (0.075);
  \draw ({\xin-1.6},\yvss) to [short,o-] ({\xin-1.3},\yvss);
  \node[anchor=east, scale=0.9] at ({\xin-1.75},\yvss) {VSS};

  % ---- the diode connection: xo1's gate up to its own drain ----
  \draw ({\xin-0.5},0) -- ({\xin-0.5},1.6) -- ({\xin+1.1},1.6);
  \fill ({\xin+1.1},1.6) circle (0.075);
  \fill ({\xin-0.5},0) circle (0.075);

  % ---- IBPS_5U arrives on the drain of the input instance ----
  \draw (XO1_d) -- ({\xin+1.1},\yout) to [short,-o] (-2.0,\yout);
  \node[anchor=east, scale=0.9] at (-2.15,\yout) {IBPS\_5U};

  % ---- both output drains leave together as IBNS_20U ----
  \draw (XO0_d) -- ({\xoa+1.1},\yout) -- ({\xob+1.1},\yout);
  \draw (XO2_d) -- ({\xob+1.1},\yout);
  \fill ({\xob+1.1},\yout) circle (0.075);
  \draw ({\xob+1.1},\yout) to [short,-o] (9.4,\yout);
  \node[anchor=west, scale=0.9] at (9.55,\yout) {IBNS\_20U};

  % ---- what the picture is for ----
  \node[anchor=north, align=center, scale=0.9, gray!50!black] at (3.7,{\yvss-0.5})
    {One device in, four devices out, so the mirror multiplies by four:\\
     5 $\mu$A into IBPS\_5U leaves as 20 $\mu$A out of IBNS\_20U};

\end{circuitikz}
\end{document}
